(1) 行星静止系的比机械能\(\varepsilon=v_\infty^2/2>0\)，比角动量\(h=bv_\infty\)。由开普勒双曲轨道\(e=\sqrt{1+2\varepsilon h^2/\mu_p^2}\)得\(e=\sqrt{1+b^2v_\infty^4/\mu_p^2}>1\)。半通径\(p=h^2/\mu_p\)，近心距离\(r_p=p/(1+e)=\{\sqrt{\mu_p^2+b^2v_\infty^4}-\mu_p\}/v_\infty^2\)，亦满足\(\frac12v_p^2-\mu_p/r_p=\frac12v_\infty^2\)。
(2) 在轨道方程\(r=p/[1+e\cos f]\)中，无穷远真近点角满足\(\cos f_\infty=-1/e\)。两条速度渐近线之间的散射偏角为\(\delta=2f_\infty-\pi=2\arcsin(1/e)=2\arctan[\mu_p/(bv_\infty^2)]\)。等价地用\(\cot(\delta/2)=bv_\infty^2/\mu_p\)，可直接看出引力越强或掠过越近，偏角越大。
(3) 行星系能量守恒保证\(|\boldsymbol u_{\rm in}|=|\boldsymbol u_{\rm out}|=v_\infty\)，且两者夹角为\(\delta\)。恒星系速度\(\boldsymbol v_{\rm in,out}=\boldsymbol V_p+\boldsymbol u_{\rm in,out}\)，于是单位测试质量的比动能改变\(\Delta K=\frac12(|\boldsymbol v_{\rm out}|^2-|\boldsymbol v_{\rm in}|^2)=\boldsymbol V_p\cdot(\boldsymbol u_{\rm out}-\boldsymbol u_{\rm in})\)。因\(|\boldsymbol u_{\rm out}-\boldsymbol u_{\rm in}|=2v_\infty\sin(\delta/2)\)，有\(|\Delta K|\le2|\boldsymbol V_p|v_\infty\sin(\delta/2)\)；增能或减能取决于遭遇几何方位。
(4) 若\(\mu_p/(bv_\infty^2)\ll1\)，\(\delta\simeq2\mu_p/(bv_\infty^2)\)，符合弱散射近似；若\(b\to0^+\)，形式上\(\delta\to\pi\)、\(r_p\to0\)，但实际有限行星半径和碰撞会在到达该数学极限前发生。恒星重力导致行星速度随时间变化，拼接圆锥近似仅在近行星遭遇时间远短于行星公转尺度时有效。
